\documentclass{article}
\usepackage{amsmath,amssymb,hyperref}
\title{How much union risk can low-order joint diagnostics hide?}
\begin{document}
\maketitle

\paragraph{Result.}
Even the exact population law of each pair of exposures need not
determine the mission-wide probability of any failure. Among
\(n\ge 3\) possible exposures, the largest gap between fixed laws
with the same labeled pair reports is
\[
  \Delta_{n,2}=
  \frac{\lfloor (n-1)^2/4\rfloor}{n+\lfloor (n-1)^2/4\rfloor}.
\]
This is an exact worst case over compatible report values, not the
interval for every report. The rule below computes the worst case
for any diagnostic order \(k\); several edge cases have exact formulas.
A low-order report that includes a justified support restriction can
be much more informative.

\paragraph{Model.}
Before admission, a single unknown fixed law governs a binary
vector \(X\) of potential failure indicators. Provided a site is
activated, its exposure yields the corresponding indicator;
other exposures do not change it in this model. There is no
resampling after a repair or deadline amendment. The observer
has no realized outcome when choosing admission. For each named
subset of at most \(k\) coordinates, the devices report its
exact marginal law. These are population constraints with no
sampling error, not a pile of unrelated empirical summaries.
The outcome is any failure among the activated sites during the
one mission horizon. The first calculation considers all \(n\)
sites active. The uncertainty set is all joint laws agreeing with
the reports and the declared support. For the worst-case gap
the support may include all binary vectors; a known restriction
on attainable counts is considered separately.

\paragraph{Reduction and sharp finite formula.}
When maximizing the gap over two laws with the same reports, apply
the same permutation averaging to both. This preserves the all-site
event and equality of the reports, though it may change their common
named values. Thus exchangeable laws suffice to attain the global
worst case. With \(Y=\sum_i X_i\), such a law is exactly a law
of \(Y\), by uniform selection within each count. For a subset of
size \(j\), its probability that all entries are one is
\(\mathbb{E}\binom Yj/\binom nj\). All lower-order assignments
are determined by these moments, using inclusion--exclusion.
For a \emph{given} symmetric report the permutation average
preserves the report itself, so the same factored model gives
exact endpoints. It need not do that for a nonsymmetric report.

Write a signed difference \(\delta\) of the two count laws.
Its moments of degrees \(0,\ldots,k\) vanish, and each sign
has mass at most one. Conversely any such difference can be
realized by adding the same leftover probability to both signs.
Orientation can be chosen so \(\delta_0\ge 0\), and the goal
is its maximum. If the optimum is positive, the sign masses
may be taken to be one. Pick an optimum with minimal support.
If it uses more than \(k+2\) counts, there is a nonzero
perturbation on that support preserving the \(k+1\) moment
equations and the mass on the positive sign. For small movement
in either direction no signs change. At an optimum the
perturbation cannot change \(\delta_0\), so move to the first
coordinate reaching zero. This contradicts minimal support.
On at most \(k+1\) counts the Vandermonde matrix has trivial
kernel. A nonzero optimum therefore uses exactly \(k+2\),
including count zero.

For \(S=\{0=x_0<x_1<\cdots<x_{k+1}\}\), put
\(\lambda_i=\prod_{j\ne i}(x_i-x_j)^{-1}\). The leading
coefficient of the Lagrange interpolant of a polynomial of degree
at most \(k\) shows that this vector annihilates all tested
moments. It spans the nullspace on \(S\). The sharp gap is
\[
 \Delta_{n,k}=\max_{S\subseteq\{0,\ldots,n\},~x_0=0,~|S|=k+2}
 \frac{2|\lambda_0|}{\sum_i|\lambda_i|}.
\]
The attaining probability laws are the two normalized sign
parts, with count zero in the low-union-risk law. If a declared
count support contains zero, the same formula restricts
\(S\) to that support. With at most \(k+1\) possible counts it
is zero by Vandermonde; if zero is impossible, the all-active
union risk is already one. In particular a genuine bound of
\(k\) simultaneous potential failures determines the risk
from order \(k\), and cannot be replaced by small marginals.

\paragraph{Closed cases.}
With one reported order, the optimum uses \(0,1,n\), giving
\(1-1/n\) by the same weight ratio with two nonzero counts.
For pair reports the four points \(0<a<b<c\) with zero
adjoined give
\[
 \frac{1}{1+ac/((b-a)(c-b))}.
\]
For fixed \(b\), choose \(a=1\) and \(c=n\) to maximize
\((b-a)(c-b)/(ac)\); choose \(b\) nearest the midpoint.
This yields the stated formula. With every proper order,
only the full set of counts is possible and binomial
alternating weights give \(\Delta_{n,n-1}=2^{1-n}\).

With order \(n-2\), omit one count \(j\in\{1,\ldots,n\}\).
The score is \(2j/\sum_t\binom nt|t-j|\). To maximize it,
use the denominator's piecewise linear interpolation in \(j\).
The derivative of the ratio has the sign of
\(\sum_{t>j}t\binom nt-\sum_{t<j}t\binom nt\) away from
integer cuts. Since \(t\binom nt=n\binom{n-1}{t-1}\),
the sign changes at the center of that shifted binomial
mass. A maximizer is \(j=\lceil n/2\rceil\); for even \(n\)
there is a flat interval through the next integer. Pairing
\(t\) and \(n-t\) in the denominator and telescoping
\((n-2t)\binom nt=n[\binom{n-1}t-\binom{n-1}{t-1}]\)
gives \(j\binom n{\lfloor n/2\rfloor}\).
Thus \(\Delta_{n,n-2}=2/\binom n{\lfloor n/2\rfloor}\).

\paragraph{Admission with an old commitment.}
For \(n=31,k=2\), the attaining count laws put masses
\((225/256,31/256)\) at \((0,16)\), or masses
\((31/32,1/32)\) at \((1,31)\). Let the known probability of
the exposure window being active be \(1/10\); outside it none
can fail. Within it the fixed law is unknown and the window
status is not observed before admission. The first two
exposures are already irrevocable and give the same risk
\(r_0=3/320\) under either law. Activating the other 29 as a
single optional bundle would make the all-active risk
\(31/2560\) or \(1/10\). This latter endpoint is feasible
and is the largest possible under the known exposure window
and common report, whether the uncertainty set has only
these two alternatives or all matching laws on that window.

Keep the original ex-ante mission cap \(1/20\). A coin
independent of the unknown realized vector admits the optional
bundle with probability \(q\). Under each fixed law the actual
risk is \(r_0+q(r_{\mathrm{all}}-r_0)\), not a reset cap and
not two independent law choices. The exact robust admission
maximum is \(13/29\); a controller informed of the low-risk
law could admit every bundle. The high-risk law still needs
the same cap. This is no guarantee conditional on knowing
the window is active. The matched planner with all named
pair reports and the same powers must use the same bound.

\paragraph{Comparators and checks.}
The binomial-moment formulation uses
\href{https://doi.org/10.1287/opre.38.2.227}{Pr\'ekopa's
sharp union bounds} as the statistical baseline. The
report-uniform polynomial certificates are also related to
\href{https://arxiv.org/abs/0803.4516}{dual polynomials for OR}.
The fixed stock of a single original risk cap is in the spirit of
\href{https://groups.csail.mit.edu/mers/papers/Ono-Williams_AAMAS10.pdf}{Ono and Williams};
the low-order experiment here isolates the effect of a single
coupled law, rather than running their trajectory optimizer or the
\href{https://arxiv.org/abs/1811.00215}{coupled robust MDP}.

The replay enumerates all diagnostic orders for \(n\le13\),
and checks 333 additional closed-form cases. Each lower witness
is checked exactly and each upper bound is an independent
rational polynomial uniform certificate on the count grid.
A second solver retains all labeled states, with 28 numerical
matches to the exact factored intervals on symmetric reports.
A third coordinate known identically zero gives a control: the
named joint report pins risk at \(3/4\), while aggregating its
moments alone allows up to \(5/6\). Lean checks the exact
admission arithmetic. This is an order-of-diagnostic result,
not a bit bound, evidence of field dependence, or R4 acceptance.
\end{document}
